\documentclass{article}

% Four-page, double-blind ML-for-Systems workshop extended abstract.
\usepackage[dblblindworkshop]{neurips_2026}
\workshoptitle{ML for Systems Workshop}
\usepackage[utf8]{inputenc}
\usepackage[T1]{fontenc}
\usepackage{hyperref}
\usepackage{url}
\usepackage{booktabs}
\usepackage{graphicx}
\usepackage{amsmath}
\usepackage{amsfonts}
\usepackage{microtype}
\usepackage{xcolor}

% This is a planning document that compiles in the conference template. Replace
% blue text as evidence becomes available; do not submit with blueprint notes.
\definecolor{blueprint}{RGB}{25,82,145}
\newcommand{\todo}[1]{\textcolor{blueprint}{[TODO: #1]}}
\newcommand{\pilot}[1]{\textcolor{blueprint}{[PILOT: #1]}}
\newcommand{\bytes}[1]{\lvert #1\rvert_{\mathrm{UTF8}}}

\title{When Tokens Expand: End-to-End Adversarial Evaluation of LLM Compression}

% Keep the submission anonymous. Add authors only in the camera-ready version.
\author{Anonymous Authors}

\begin{document}
\maketitle

\begin{abstract}
Language-model compressors replace a hand-designed predictor with learned
next-token probabilities, but their system behavior depends on more than model
accuracy. Tokenization is a variable-rate interface between source bytes and
model symbols; probability quantization, arithmetic coding, dictionary
transmission, and framing add further costs. We present a principled methodology
for adversarially evaluating this complete pipeline. Our attacks preserve exact
token and byte round trips while optimizing either surprisal per decoded byte or
realized serialized size. Layer-by-layer accounting distinguishes model failure
from coder and container overhead, and matched conventional codecs test whether
failures are specific to learned compression. A Qwen2.5-0.5B evaluation finds that three canonical adversarial sequences
require 5.73 times the ideal bits per token of matched text8. Post-hoc
occurring-token masks reduce adversarial cost by 35.2\%, but it remains
3.81 times the masked text8 cost. These results motivate source-normalized robustness metrics and a guarded codec that falls back
to a conventional compressor when the learned path expands its input.
\end{abstract}

% -----------------------------------------------------------------------------
% BLUEPRINT AT A GLANCE (delete before submission)
% Target: at most four content pages, including figures and tables. Suggested
% budget: motivation 0.75; formulation and attacks 1.0; evaluation 0.75; results
% 0.9; systems implications, limitations, and conclusion 0.6. Do not assume
% references or supplementary material are excluded from the organizers four-page
% limit; confirm their counting rule before submission.
%
% One-sentence thesis:
% Tokenization is a variable-rate front end to predictive coding, so robust
% evaluation must optimize and report code length per decoded source byte, not
% merely next-token surprisal.
%
% Claim discipline:
% (C1) Formal: token surprisal and source-normalized compression ratio are
%      different objectives whenever decoded token lengths vary.
% (C2) Method: canonical, byte-aware greedy and beam attacks directly optimize
%      entropy-floor or realized arithmetic-code expansion.
% (C3) Empirical: use only after the full multi-model/seed matrix supports it.
%      The current repository establishes a Qwen2.5-0.5B case study, not broad
%      tokenizer or model universality.
% -----------------------------------------------------------------------------

\section{A learned predictor is only one system component}

The systems problem is not merely to replace a hand-crafted predictor with an
LLM. The operational object is a pipeline from UTF-8 bytes through tokenization,
model inference, frequency quantization, arithmetic coding, dictionary transfer,
and container framing. We ask whether untrusted inputs can make that complete
system expand data, which layer causes the expansion, and how the system should
respond. Our contribution is an end-to-end methodology with explicit byte-level
contracts, system-aware attacks, cost attribution, and a bounded fallback path.


In this pipeline, tokenizer symbols represent unequal numbers of bytes and may
decode context-dependently; finite frequency tables distort probabilities; a
restricted vocabulary must be transmitted; and headers dominate small inputs.
Thus the least likely token need not maximize compressed bits per source byte,
and perplexity alone cannot characterize end-to-end robustness. We call
tokenization a \emph{compression interface}, not a self-contained compressor:
the vocabulary and token-ID code must still be shared or transmitted.

\paragraph{Contributions.}
We make three systems contributions: (i) source-normalized metrics that separate
model NLL, coder payload, dictionary transmission, and serialized size; (ii)
canonical greedy and beam attacks that optimize decoded-byte or realized-code
objectives; and (iii) a matched evaluation with random, natural-text, and
conventional-codec controls. \todo{Add the supported multi-model conclusion.}

\paragraph{Scope.}
We do not claim that tokenization alone is a self-contained compressor, that
approximate beam search finds a global worst case, or that sequence NLL is a
universal lower bound for an individual string.

\section{Tokenization as a compression interface}
\label{sec:formulation}

Let $s$ be UTF-8 text, $T(s)=x_{1:N}$ its tokens, $D(x)$ its decoding,
$h_t=x_{1:t-1}$, and $p_\theta(v\mid h_t)$ the model distribution. If the
codec retains dictionary $A\subseteq V$, it renormalizes that distribution to
$q_\theta(v\mid h_t,A)$. The realized sequence NLL is
\begin{equation}
  L_A(x)=-\sum_{t=2}^{N}\log_2q_\theta(x_t\mid h_t,A).
\end{equation}
This is sequence cross-entropy, not Shannon entropy of the predictive
distribution.

\subsection{Source-normalized cost}

Define the source-normalized ideal rate
\begin{equation}
  R_{\mathrm{NLL}}(x;A)=\frac{L_A(x)}{8\bytes{D(x)}}.
  \label{eq:ideal-rate}
\end{equation}
For a realized arithmetic payload $C_{\mathrm{AC}}$, fixed or transmitted
dictionary cost $C_A$, and framing/header cost $C_H$, report
\begin{align}
  R_{\mathrm{payload}} &= \frac{C_{\mathrm{AC}}}{8\bytes{D(x)}}, \\
  R_{\mathrm{serialized}} &=
  \frac{C_{\mathrm{AC}}+C_A+C_H}{8\bytes{D(x)}}.
  \label{eq:serialized-rate}
\end{align}
Values above one denote expansion. The gap between
$R_{\mathrm{NLL}}$ and $R_{\mathrm{payload}}$ measures coder/quantization cost;
the remaining gap to $R_{\mathrm{serialized}}$ measures metadata and framing.

\paragraph{Objective mismatch.}
For candidate $v$, let $\Delta_b(v;h)=\bytes{D(h\mathbin\Vert v)}-
\bytes{D(h)}$. Ranking by $-\log p(v\mid h)$ need not match ranking by
$-\log p(v\mid h)/\Delta_b(v;h)$: a 12-bit-surprise token adding four bytes
loses to a 6-bit-surprise token adding one byte under the latter objective.

\subsection{Canonicality and losslessness}

Tokenizers need not satisfy $T(D(x))=x$ for arbitrary token-ID sequences. An
attack sequence is valid only if every accepted prefix obeys
\begin{equation}
  T(D(x_{1:t}))=x_{1:t}.
  \label{eq:canonical}
\end{equation}
The implementation examines candidates in objective order and accepts the first
canonical extension. By induction, the completed sequence is the canonical token
representation of the bytes supplied to every compressor. Final validation must
also decode the serialized arithmetic-coded artifact and compare both token IDs
and source bytes exactly.

\section{Threat model and attacks}
\label{sec:attacks}

\paragraph{Adversary.}
Use a white-box adversary with access to tokenizer, logits, and coding algorithm.
The adversary chooses a length-$N$ canonical token sequence after a fixed seed,
subject to a candidate alphabet $G$. Evaluate three alphabets separately:
full non-special vocabulary, printable ASCII bytes, and canonical one-byte ASCII.
The last isolates predictive failure from token-length variation; the broader
alphabets expose the tokenizer interface itself.

\paragraph{Attacks.}
We compare seeded uniform valid tokens; greedy minimum probability; greedy
maximum surprisal per added byte; beam search over cumulative NLL/source bit; and
beam search that clones arithmetic-coder state and ranks finalized size. Byte
increments decode the complete prefix before and after extension. We test the
full non-special vocabulary, printable ASCII, and canonical one-byte ASCII. Beam
search is approximate unless its branch factor covers all candidates; changing
beam ancestry also requires disabling the KV cache.

\paragraph{Global mask.}
We construct each adversarial sequence once under the full distribution, then
replay identical IDs and histories under its occurring-token mask. This isolates
renormalization and is not a mask-adaptive attack.

\section{Experimental evaluation}
\label{sec:design}

\paragraph{Setup.}
We evaluate Qwen2.5-0.5B with its 151,643-token vocabulary. The adversary creates
three 1,000-token sequences from starting token IDs 785, 32, and 641 by repeatedly
choosing the lowest-probability canonical extension. Every completed sequence
satisfies $T(D(x))=x$ and decodes to 6,988, 6,812, and 6,943 UTF-8 bytes. We then
replay the identical token IDs and contexts using the set of tokens occurring in
each completed sequence; the three post-hoc dictionaries contain 592, 606, and
587 tokens. Thus masking changes only probability normalization, not generation.

The natural-text control is the first 100,000 text8 tokens, processed by the same
model with context length 1,000, retained context 100, KV caching, and arithmetic
coding. It uses 16 independently seeded batches, so 99,984 tokens are predicted.
We report mean adversarial surprisal with sample standard deviation across the
three starts. Text8 is one deterministic matched run per dictionary policy and
therefore has no error bar.

\begin{figure}[t]
  \centering
  \includegraphics[width=0.98\linewidth]{../artifacts/figures/neurips/typical_vs_adversarial.pdf}
  \caption{Notebook results for matched Qwen2.5-0.5B runs. Left: ideal
  sequence NLL per predicted token; adversarial bars show mean $\pm$ sample SD
  across three starts. Right: text8 NLL, realized arithmetic payload, and total
  including the dictionary bitmap. The adversarial result is idealized NLL, not
  a realized arithmetic payload.}
  \label{fig:typical-worst}
\end{figure}

\paragraph{Worst case versus typical text.}
Figure~\ref{fig:typical-worst} shows that full-vocabulary adversarial inputs
require $28.409\pm0.077$ ideal bits per predicted token, compared with 4.954 for
text8---a $5.734\times$ increase. The adversarial sequences use on average 595
distinct tokens. Restricting each fixed sequence to its occurring-token
dictionary lowers its cost to $18.402\pm0.833$ bits/token, a 35.2\% reduction,
but it remains $3.806\times$ the matched text8 cost of 4.835 bits/token. Identical
full-model scores under both policies verify that the token IDs and contexts did
not change.

For text8, the full-vocabulary arithmetic payload is 6.136 bits per predicted
token and the total is 6.138 bits per source token. The occurring-token policy
reduces these to 4.860 and 5.724 bits, respectively, despite transmitting a
10,818-byte bitmap. Relative to raw UTF-8 size, the totals are 14.405\% and
13.434\%, corresponding to $6.942\times$ and $7.444\times$ compression. At
100,000 tokens, the occurring-token bitmap is 2.03\% of raw size, so its coding
gain exceeds its transmission cost.

\paragraph{Source-byte interpretation.}
The adversarial sequences have more UTF-8 bytes per token than text8. Relative to
their own measured, round-tripping bytes, their mean model-induced floors are
51.346\% for the full vocabulary and 33.284\% after masking. These are floors
for codes constrained to the stated model distributions; they exclude finite-
coder and dictionary overhead and are not universal bounds for an individual
string. This distinction is precisely why both token and source-byte views are
needed.

\paragraph{Coder-aware validation.}
A preliminary 32-token one-byte-ASCII run provides a sanity check: token- and
byte-greedy attacks coincide at 135.48\% NLL and 143.36\% payload, while an NLL
beam reaches 136.70\% and 144.14\%. Because beam width is two and fixed overhead
dominates, we use this only to validate objectives and accounting.

\section{Systems implications and limitations}
\label{sec:limitations}

The systems lesson is that learned components need contracts at conventional
boundaries. Source-normalized accounting identifies whether fragility comes from
the tokenizer, predictor, quantizer, coder, dictionary, or container. For
untrusted inputs, a guarded codec can store the smaller of learned and fallback
encodings plus a selector bit, bounding size by the conventional codec plus one
bit; its latency and throughput cost remains to be measured.

Our evidence is limited to one model, synthetic white-box inputs, three starts,
finite context, approximate search, and one coder/container. It does not establish
global worst cases or generalize to multilingual text or code. The attack could
enable storage or compute amplification; practical mitigations include size caps,
early-abort expansion detection, bounded inference, and fallback codecs.

\section{Conclusion}

Robust learned compression is an end-to-end systems problem. Canonical byte-aware attacks and explicit cost attribution show that the tested
adversarial sequences are 3.81--5.73 times harder than matched text8, while
guarded fallback offers a principled path from diagnosis to robust operation.

\begin{ack}
\todo{Camera-ready only: funding and competing-interest disclosures. The supplied
style hides this environment in anonymous submission mode.}
\end{ack}

\section*{References}
\todo{Add a BibTeX database of verified primary sources and use one consistent
natbib citation style. References do not count toward the content-page limit.}

% Supplementary planning notes (do not count toward the four-page narrative):
% - Prove the token/byte objective mismatch and give attack pseudocode.
% - Record configurations, revisions, commands, seeds, hardware, compute, and
%   licenses; report full per-run tables and search-budget ablations.
% - Document exact seed-token, quantization, bitmap, payload, and header costs.
% - Repository map: src/adversarial.py; src/compression_attacks.py;
%   scripts/run_compression_attacks.py; evaluation/notebooks/neurips/
%   adversarial_worst_case.ipynb; and artifacts/runs/adversarial/
%   qwen_05b_ascii_bytes_n1000.
% - Confirm whether references and the NeurIPS checklist are outside the
%   organizers four-page limit. Keep the checklist unless organizers explicitly
%   say their workshop does not require it.

\newpage
\input{checklist.tex}

\end{document}
